\section{The Selection Argument}

The whole weight of the theory rests on one claim, that the substrate is selected and not chosen, forced by a short list of demands rather than fitted to the answer. This section makes that case. It has two parts, a survey that shows what is common and what is rare across the whole family of honeycombs, and a ladder of demands that, together, leave exactly one survivor.

The survey runs a single battery across every regular hyperbolic tessellation, all forty-two that can be built, measuring each the same way so the numbers are directly comparable. The first result is that matter is cheap. A propagating fermion, the Kahler-Dirac field of the comparison chapter, lives and moves on all forty-two, its clean return probability falling with size on every one. Hyperbolic geometry as such is enough to carry a moving relativistic particle. The second result is that the spinor sites are dear. The native spinor structure, the directions that carry spin without a detour, appears on only seven of the forty-two, exactly the honeycombs faceted by the $24$-cell, and these are dimension-gated, appearing first in four dimensions. Matter is everywhere, but the structure that carries the gauge symmetry and the generations is rare and four-dimensional.

That rarity is what the selection exploits. The argument that the substrate must be $\{3,4,3,4\}$ has two independent teeth, and only one honeycomb survives both.

\begin{table}[ht]
\centering
\begin{tabular}{@{}lp{0.6\textwidth}@{}}
\toprule
demand & what it forces \\
\midrule
three-dimensional space & physical space is the cusp, one dimension below the bulk, so observing 3D space forces a 4D bulk \\
crystallographic & the honeycomb symbol may contain only $3$s and $4$s, ruling out every member with a $5$ \\
spinor in the sites & the directions must carry spin natively, the $24$-cell faceting \\
hyperbolic & the space must grow without bound and without a preferred direction \\
\bottomrule
\end{tabular}
\caption{The four demands. Exactly one regular honeycomb meets all four.}
\label{tab:selection}
\end{table}

The first tooth is dimension. Because physical space is the flat cusp and the cusp is always one dimension below the bulk, the three-dimensional space we observe forces a four-dimensional bulk, no more and no less. The second tooth is the crystallographic spinor. The compact regular honeycombs that exist in two, three, and four dimensions all carry a $5$ in their symbol, the icosahedral number, and a $5$ both breaks the crystallographic restriction and, as the comparison chapter showed, blocks the native spinor. Demanding both crystallography, only $3$s and $4$s, and a spinor in the sites eliminates every compact regular honeycomb and leaves the paracompact ones. Within four dimensions, the one honeycomb that is at once crystallographic, hyperbolic, spinor-carrying, and equipped with a flat three-dimensional cusp is $\{3,4,3,4\}$. Its flat twin $\{3,4,3,3\}$ has the same spinor sites but no curvature, so it fails the last demand, and curvature is exactly what separates the committed substrate from that twin.

\begin{result}[The unique substrate]
$\{3,4,3,4\}$ is the only regular honeycomb that is simultaneously crystallographic, hyperbolic, spinor-carrying in its sites, and three-dimensional in its cusp. The substrate is therefore selected by four demands, not chosen to fit, which is the strongest form the argument takes (depth: derived and measured).
\end{result}

This does not reduce the arbitrariness to zero, and the paper does not claim it does. The four demands are themselves choices, motivated but not proven inevitable, and a reader is free to weigh them. What the selection argument does is turn a free parameter, which crystal, into a forced consequence of a few stated requirements, which is a real reduction even if it is not a derivation from nothing. The substrate is the narrowest part of the theory, and it earns that narrowness here.
